\section{Problema del Generatore di Tensione}

Il metodo dei potenziali ai nodi si basa sull’applicazione della Legge di Kirchhoff delle correnti (LKC) a ciascun nodo del circuito, escluso il riferimento. Quando sono presenti generatori ideali di tensione si presenta una difficoltà: un generatore di tensione impone una differenza di potenziale nota tra due nodi, ma la corrente che lo attraversa è incognita e dipende dal resto del circuito. Non è quindi possibile esprimere direttamente la corrente nel ramo del generatore tramite la legge di Ohm.

Per risolvere il problema si introduce una procedura sistematica che:
\begin{itemize}
    \item mantiene i potenziali nodali come incognite principali;
    \item aggiunge una variabile per ogni generatore di tensione: la corrente che lo attraversa (indicata con \(I_E\));
    \item scrive le equazioni di KLC a ciascun nodo includendo queste correnti incognite;
    \item aggiunge un’equazione di vincolo per ogni generatore, che lega i potenziali nodali alla tensione imposta.
\end{itemize}

\subsection*{Esempio applicativo}

Applichiamo la procedura al circuito di Figura \ref{fig:circuito_nodi_1}, che contiene tre nodi indipendenti e due generatori di tensione. I potenziali nodali \(V_A, V_B, V_C\) sono stati definiti rispetto al nodo di riferimento \(D\) (massa).

\CircuitoPotenzialeNodi{3}  

Il procedimento è molto simile a quello spiegato nel capitolo precedente, con l’aggiunta del punto 3.1:

\begin{enumerate}
    \item[1.] \textcolor{gray}{\textbf{Identificare i nodi e scegliere il nodo di riferimento:} Si individua il nodo di riferimento (massa) e si assegnano potenziali incogniti agli \(N-1\) nodi rimanenti.}
    \item[2.] \textcolor{gray}{\textbf{Scrivere le equazioni nodali (LKC):} Per ogni nodo diverso dal riferimento si impone la somma algebrica delle correnti nulla.}
    \item[3.] \textcolor{gray}{\textbf{Applicare Ohm:} Utilizzare la legge di Ohm per esprimere le correnti nei resistori in funzione dei potenziali nodali.}
    \item[3.1.] \textbf{Equazioni di vincolo dei generatori di tensione:} Per ogni generatore di tensione si aggiunge una equazione che lega la differenza di potenziale tra i due nodi a cui è collegato al valore della tensione impressa; inoltre si introduce la corrente incognita del generatore nelle equazioni di KCL dei nodi interessati.
    \item[4.] \textcolor{gray}{\textbf{BONUS: Matrici!} Organizzare le equazioni in forma matriciale, identificando la matrice delle conduttanze e il vettore dei termini noti.}
\end{enumerate}

\subsection*{1. Identificare i nodi e scegliere il nodo di riferimento}

Nel circuito di Figura \ref{fig:potenziale_3} sono presenti \(N=4\) nodi. Scegliamo come riferimento il nodo \(B\) (massa). Assegniamo quindi i potenziali incogniti \(V_A\), \(V_C\) e \(V_D\) ai nodi rimanenti.

\CircuitoPotenzialeNodi{4}  % figura da definire

\hrule
\subsection*{2. Scrivere le equazioni nodali (LKC)}

Per ogni nodo diverso dal riferimento (nodi \(A\), \(C\) e \(D\)) imponiamo la somma algebrica delle correnti entranti nulla. Nelle equazioni compare anche la corrente incognita dei generatori di tensione, che indichiamo con \(I_{E_1}\) per il generatore \(E_1\) e \(I_{E_2}\) per \(E_2\). Consideriamo il verso della corrente nei rami con i generatori di tensione uscente dal morsetto positivo degli stessi.

\begin{itemize}
    \item \textcolor{red}{\textbf{Nodo A:}}  Il nodo $A$ è influenzato dal generatore di tensione $E_1$, quindi dobbiamo tenere conto della corrente che circola su quel ramo.
    \begin{equation}
        J_1 - I_{R_1} + I_{E_1} = 0
        \label{eq:nodoA_vincolo}
    \end{equation}
    \item \textcolor{blue}{\textbf{Nodo C:}} Il nodo $C$ è influenzato dal generatore di tensione $E_1$.
    \begin{equation}
        -J_1 - I_{R_1} - I_{E_1} = 0
        \label{eq:nodoC_vincolo}
    \end{equation}
    \item \textcolor{magenta}{\textbf{Nodo D:}} Il nodo $D$ è influenzato dal generatore di tensione $E_1$.
    \begin{equation}
        J_2 + I_{E_1} - I_{R_4} = 0
        \label{eq:nodoD_vincolo}
    \end{equation}
\end{itemize}

\vspace*{0.3cm}

\hrule




\subsection*{3. Applicare la legge di Ohm}

Per ogni equazione bisogna definire ogni corrente come rapporto tra resistenza complessiva sul ramo e tensione tra i nodi. Inoltre riscriviamo sempre l'equazione utilizzando il concetto di \textbf{conduttanza}.

\paragraph{Nodo A}
L'unica corrente incognita è la corrente che attraversa il resistore $R_4$, resistore che si collega ai morsetti $A$ e $D$. Perciò l'equazione diventa:
\begin{equation}
J_1 - J_4 - \frac{V_A - V_D}{R_4} = 0 \qquad \Rightarrow \qquad J_1 - J_4 - G_4\left(V_A - V_D\right) = 0
\label{eq:nodoA_Ohm}
\end{equation}

\paragraph{Nodo C}
L'unica corrente incognita è la corrente che attraversa il resistore $R_1$, resistore che si collega ai morsetti $C$ e $B$. Notiamo però che la tensione sul morsetto $B$ è pari a 0, visto che è collegato al GND. Perciò l'equazione diventa:
\begin{equation}
- J_1 + J_2 - \frac{V_C - V_B}{R_1} = 0 \qquad \Rightarrow \qquad - J_1 + J_2 - G_1\left(V_C\right) = 0
\label{eq:nodoC_Ohm}
\end{equation}


\paragraph{Nodo D}
L'unica corrente incognita è la corrente che attraversa il resistore $R_4$, resistore che si collega ai morsetti $A$ e $D$. Perciò l'equazione diventa:
\begin{equation}
- J_2 + J_3 - \frac{V_D - V_A}{R_4} = 0 \qquad \Rightarrow \qquad - J_2 + J_3 - G_4\left(V_D-V_A\right) = 0
\label{eq:nodoD_Ohm}
\end{equation}




Per una migliore comprensione, possiamo mettere a sistema le equazioni trovate e riorganizziamo i termini per esplicitare le varie incognite.
Quindi il nostro sistema si potrà riscrivere in questa maniera:

\begin{equation}
\left\{
\begin{aligned}
G_4\, & \textcolor{red}{V_A} & +\; 0\, & \textcolor{blue}{V_C} & -\; G_4\, & \textcolor{magenta}{V_D} & = & J_1 - J_4 \\
0\, & \textcolor{red}{V_A} & +\; G_1\, & \textcolor{blue}{V_C} & +\; 0\, & \textcolor{magenta}{V_D} & = & J_2 - J_1 \\
-G_4\, & \textcolor{red}{V_A} & +\; 0\, & \textcolor{blue}{V_C} & +\; G_4\, & \textcolor{magenta}{V_D} & = & -J_2 + J_3
\end{aligned}
\right.
\end{equation}








\hrule
\subsection*{3.1. Equazioni di vincolo dei generatori di tensione}

Ogni generatore di tensione impone una relazione tra i potenziali dei nodi a cui è collegato.

\begin{itemize}
    \item Generatore \(E_1\): è collegato tra il nodo \(A\) e il nodo \(B\). La tensione impressa, con la polarità indicata in figura, è:
          \begin{equation}
              V_A - V_B = E_1
              \label{eq:vincoloE1}
          \end{equation}
    \item Generatore \(E_2\): è collegato tra il nodo \(B\) e il nodo \(C\). La tensione impressa è:
          \begin{equation}
              V_B - V_C = E_2
              \label{eq:vincoloE2}
          \end{equation}
\end{itemize}

Queste due equazioni, insieme alle tre equazioni di KCL, formano un sistema di 5 equazioni nelle 5 incognite \(V_A, V_B, V_C, I_{E_1}, I_{E_2}\).

Riorganizziamo il sistema in forma canonica, portando tutti i termini incogniti a sinistra e i termini noti a destra:

\begin{equation}
    \left\{
    \begin{aligned}
        -(G_1 + G_3)\, V_A & + 0\, V_B & + G_3\, V_C & + 1\, I_{E_1} & + 0\, I_{E_2} & = -J_1 \\
        0\, V_A            & - G_2\, V_B & + G_2\, V_C & - 1\, I_{E_1} & + 1\, I_{E_2} & =  J_2 \\
        G_3\, V_A          & + G_2\, V_B & -(G_2 + G_3)\, V_C & + 0\, I_{E_1} & - 1\, I_{E_2} & =  0   \\
        1\, V_A            & -1\, V_B    & + 0\, V_C & + 0\, I_{E_1} & + 0\, I_{E_2} & =  E_1 \\
        0\, V_A            & +1\, V_B    & -1\, V_C & + 0\, I_{E_1} & + 0\, I_{E_2} & =  E_2
    \end{aligned}
    \right.
\end{equation}

\hrule
\subsection*{4. BONUS: Matrici!}

Anche in questo caso la rappresentazione matriciale semplifica la gestione del sistema.

\subsubsection*{Matrice dei coefficienti}

La matrice dei coefficienti \(\mathbf{A}\) può essere suddivisa in blocchi che riflettono la struttura del problema:

\begin{equation}
    \mathbf{A} = 
    \begin{bmatrix}
        \mathbf{G} & \mathbf{C} \\
        \mathbf{E} & \mathbf{0}
    \end{bmatrix}
\end{equation}

dove:
\begin{itemize}
    \item \(\mathbf{G}\) è la matrice delle conduttanze nodali (3×3 nel nostro esempio), costruita con la regola classica: \(G_{ii}\) somma delle conduttanze incidenti al nodo \(i\), \(G_{ij} = -G_{ij}\) (con segno negativo) per i rami resistivi tra nodi \(i\) e \(j\).
    \item \(\mathbf{C}\) (3×2) descrive come le correnti dei generatori di tensione compaiono nelle equazioni di KCL: \(C_{ik} = +1\) se la corrente \(I_{E_k}\) entra nel nodo \(i\), \(-1\) se ne esce, \(0\) altrimenti.
    \item \(\mathbf{E}\) (2×3) è la matrice dei vincoli di tensione: \(E_{kj} = +1\) se il potenziale \(V_j\) compare con segno positivo nella relazione del generatore \(k\), \(-1\) se con segno negativo, \(0\) altrimenti.
    \item Lo zero in basso a destra indica che i generatori di tensione non hanno equazioni che legano tra loro le loro correnti.
\end{itemize}

Nel nostro esempio:

\begin{equation}
    \mathbf{G} = \begin{bmatrix}
        -(G_1+G_3) & 0 & G_3 \\
        0 & -G_2 & G_2 \\
        G_3 & G_2 & -(G_2+G_3)
    \end{bmatrix},\qquad
    \mathbf{C} = \begin{bmatrix}
        1 & 0 \\
        -1 & 1 \\
        0 & -1
    \end{bmatrix},\qquad
    \mathbf{E} = \begin{bmatrix}
        1 & -1 & 0 \\
        0 & 1 & -1
    \end{bmatrix}.
\end{equation}

La matrice completa è quindi:

\begin{equation}
    \mathbf{A} = 
    \begin{bmatrix}
        -(G_1+G_3) & 0 & G_3 & 1 & 0 \\
        0 & -G_2 & G_2 & -1 & 1 \\
        G_3 & G_2 & -(G_2+G_3) & 0 & -1 \\
        1 & -1 & 0 & 0 & 0 \\
        0 & 1 & -1 & 0 & 0
    \end{bmatrix}.
\end{equation}

\subsubsection*{Vettore incognita}

Il vettore incognita contiene tutti i potenziali nodali e tutte le correnti dei generatori di tensione:

\begin{equation}
    \mathbf{x} = 
    \begin{bmatrix}
        V_A \\ V_B \\ V_C \\ I_{E_1} \\ I_{E_2}
    \end{bmatrix}.
\end{equation}

\subsubsection*{Vettore dei termini noti}

Il vettore \(\mathbf{b}\) raccoglie i contributi dei generatori indipendenti di corrente e delle tensioni impresse:

\begin{equation}
    \mathbf{b} = 
    \begin{bmatrix}
        -J_1 \\ J_2 \\ 0 \\ E_1 \\ E_2
    \end{bmatrix}.
\end{equation}

In forma compatta:

\begin{equation}
    \mathbf{A} \, \mathbf{x} = \mathbf{b}.
\end{equation}

La struttura a blocchi evidenzia come la presenza dei generatori di tensione estenda il sistema originale delle conduttanze con nuove variabili e nuove equazioni. Questo approccio permette di applicare metodi risolutivi standard (eliminazione di Gauss, fattorizzazione LU) anche in presenza di generatori di tensione ideali, senza dover ricorrere al concetto di supernodo.